\section{Introduction}\label{sec:intro}

\begin{figure}[H]
  \centering
  \includegraphics[width=\columnwidth]{\ksFigure{bridge_conv0.pdf}}
  \caption{The exact-to-heuristic evidence ladder on Conv\_Case0, decomposed by
  $\mathrm{Tr}=\mathrm{Vol}+\mathrm{Dt}$. The scalable production solver
  improves the official artifact by $9.1\%$; exploratory order repair adds
  $1.9\%$; and the weighted-gap planner attains its lower bound under fixed
  dependency-frontier, certifying 57,408 bytes.}
  \Description{Horizontal bars compare official, scalable, repair, and exact
  plans, split into backed volume, generated volume, and generated surcharge.}
  \label{fig:bridge}
\end{figure}

NPU compilers lower each kernel into thousands of micro-operations executing on
heterogeneous pipelines and exchanging data through software-managed on-chip
memories. Scheduling commits simultaneously to a legal topological order of the
micro-operation DAG, contiguous physical addresses for resident buffers, and a
spill plan when capacity is insufficient. Allocation, use, and release timing
determines both how many bytes leave the chip and which values are evicted;
logical live peaks alone capture neither fragmentation, next-use distance, nor
asymmetric spill cost.

The benchmark charges one transfer for input-backed buffers and two for
generated, unbacked buffers. If $\mathrm{Vol}$ is total spilled volume and
$\mathrm{Dt}$ its generated part, then $\mathrm{Tr}=\mathrm{Vol}+\mathrm{Dt}$.
Traffic changes must therefore be explained through both total volume and the
generated surcharge. Section~\ref{sec:formulation} gives the formal definition
and tight bounds; Supplementary Sec.~2.1 gives the complete proof.

We separate topological-order construction from physical spill planning and use
a machine-checkable exact planner to quantify scalable-heuristic headroom
(Figure~\ref{fig:bridge}). The production solver uses a
\emph{dependency-frontier} order to prevent starvation of ready predecessor
groups, then selects by the true P2/P3 objective among a small set of structural
orders, placement/eviction policies, and reload windows. A fixed-order weighted
residency-gap oracle supplies a traffic lower bound; validated contiguous
packing that attains the bound yields a fixed-order traffic certificate. Offline
repair and exact planning are reported as research tools, not default stages.
The complete solution process appears in Supplementary Sec.~3.

Our contributions are: (1) a scalable dependency-frontier scheduler integrated
with a true-objective spill-planning portfolio; (2) a weighted residency-gap
lower bound that upgrades to a fixed-order certificate through a validated
contiguous artifact; and (3) an audit-oriented evaluation organized around
$\mathrm{Tr}=\mathrm{Vol}+\mathrm{Dt}$ and a strict separation of production,
repair case studies, and the exact oracle. On six public instances, production
has five P2 wins and one tie and five P3 time wins with one loss. The Conv\_Case0
oracle certifies 57,408 bytes and exposes a $29.6\%$ fixed-order difference from
the legacy P1 order. Nonuniform repair and internal synthetic results are
additional bounded evidence only; see Section~\ref{sec:experiments} and
Supplementary Sec.~4.

The source code and reproduction guide are available online.\footnote{\url{https://github.com/VennIntelligence/2025A-kernel-sched}}
